% !PRINTED 259
\weberpage{259}

%<<CONT
bestimmt werden, während die $\alpha$ neue Coëfficienten sind, die in
einer leicht zu übersehenden Weise aus den $\beta$ abgeleitet werden,
auf deren Bildung es uns hier nicht weiter ankommt.

Nun lässt sich nach der Voraussetzung die Function
$\varphi(x_1, x_2 \ldots x_m)$ durch $z_1, z_2 \ldots z_k$, also auch durch $y_1, y_2 \ldots y_k$
allein ausdrücken. Bezeichnen wir diesen Ausdruck mit
$\Phi(y_1, y_2 \ldots y_k)$, so haben wir [durch Vermittelung von (3)] die
Identität
\weq{4}{\varphi(x_1, x_2 \ldots x_k, x_{k+1} \ldots x_m) = \Phi(y_1, y_2 \ldots y_k),}
und wenn wir darin $x_{k+1}, x_{k+2}, \ldots x_m = 0$ setzen,
\weq{5}{\Phi(x_1, x_2 \ldots x_k) = \varphi(x_1, x_2 \ldots x_k, 0, 0 \ldots 0),}
wodurch, da es auf die Bezeichnung der Variablen nicht ankommt,
$\Phi$ vollständig bestimmt ist; wir können daher setzen
\begin{align}
\Phi(y_1, y_2 \ldots y_k) &= \varphi(y_1, y_2 \ldots y_k, 0, 0 \ldots 0) \tag{6}\\
&= \varphi(x_1, x_2 \ldots x_m);
\end{align}
$\Phi$ entsteht also aus $\varphi$ dadurch, dass man die $m-k$ letzten
Variablen (bei der hier gewählten Anordnung) gleich Null setzt:
\weq{7}{\Phi(x_1, \ldots x_k) = \sum_{1,k}^{r,s} a_{r,s} x_r x_s.}

Die Determinante von $\Phi$ ist eine der Haupt-Unterdeterminanten von $R$, und zwar erhält man sie, indem man in $R$ die
$m-k$ letzten Zeilen und Colonnen wegstreicht.

Nehmen wir an, dass $\varphi$ sich nicht durch noch weniger als
$k$ Variable ausdrücken lässt, dass also $k = m - \varrho$ sei, so kann
die Determinante dieser Function $\Phi$ nicht verschwinden.

Hieraus und aus dem zuvor Bewiesenen ergiebt sich nun
der Satz:

\gesperrt{II.} \gesperrt{Wenn alle Haupt-Unterdeterminanten von $R$ von
mehr als $k$ Reihen verschwinden, während unter
den Haupt-Unterdeterminanten von $k$ Reihen
wenigstens eine von Null verschieden ist, so ist}
\[
\varrho = m - k.
\]
Denn aus (7) folgt, dass, wenn $\varrho = m - k$ ist, wenigstens
eine $k$-reihige Haupt-Unterdeterminante von Null verschieden
sein muss, und aus (1), dass, wenn auch noch eine Haupt-Unterdeterminante von mehr als $k$ Reihen von Null verschieden ist,
\[
\varrho < m - k
\]
ist.
